%=====================================================================
\chapter{Informed Search and Heuristics}\label{ch:informed}
%=====================================================================
\locator{2}{Part II --- Problem Solving by Search}

\begin{outcomes}
By the end of this chapter you will be able to:
\begin{tight}
  \item \textbf{Define} a heuristic function and \textbf{distinguish} greedy
        best-first search from $A^{*}$ by the term each one minimises.
  \item \textbf{State} admissibility and consistency precisely, and
        \textbf{explain} which one $A^{*}$ needs for tree search and which for
        graph search.
  \item \textbf{Prove} that $A^{*}$ with an admissible heuristic returns an optimal
        solution, and \textbf{identify} the step of the proof that fails when the
        heuristic overestimates.
  \item \textbf{Invent} an admissible heuristic for a new problem by relaxing a
        constraint, and \textbf{compare} two heuristics by dominance.
  \item \textbf{Measure} the trade-off between nodes expanded and solution quality,
        and \textbf{choose} greedy or $A^{*}$ with a justification.
\end{tight}
\end{outcomes}

\begin{prereq}
\chref{uninformed} (uniform-cost search, the explored set, goal-testing on
popping, and why every uninformed strategy is exponential) and \chref{python} (the
\texttt{heapq} frontier with a tie-breaking counter). $A^{*}$ is uniform-cost
search with one term added, so the previous chapter's code is the starting point.
\end{prereq}

\begin{keyterms}
heuristic function $\bullet$ greedy best-first search $\bullet$ $A^{*}$ $\bullet$
evaluation function $\bullet$ path cost $g$ $\bullet$ heuristic estimate $h$
$\bullet$ admissibility $\bullet$ consistency $\bullet$ monotonicity $\bullet$
dominance $\bullet$ relaxed problem $\bullet$ critical path $\bullet$ effective
branching factor $\bullet$ optimally efficient
\end{keyterms}

\section{A Sprint Where the Cost Matters}

\chref{uninformed} ended on an awkward note: every complete plan for the release
cost the same twenty-seven days, because every task was mandatory. Cost-based
search had nothing to choose between. Real sprints are not like that --- there is
almost always more than one way to get a thing done.

So this chapter uses a slightly larger network with one genuine choice in it. The
team must get a release out, and the testing can be done either way:

\begin{itemize}[leftmargin=2.2em, itemsep=2pt, topsep=4pt]
  \item \textbf{Manual testing} --- six days, needs only the API and authentication
        finished.
  \item \textbf{Automated testing} --- two days, but it needs a test harness built
        first, which is three days. Five days in total, and the harness needs the
        schema.
\end{itemize}

Automation wins by one day, and the whole point of this chapter is that finding out
\emph{which} route is cheaper is a search problem, and that a good heuristic finds
it while doing a third less work. \dsref{sprint-net} is the network.

\dsfig{%
\begin{tikzpicture}[font=\scriptsize, >={Stealth[length=1.8mm]},
  t/.style={rounded corners=3pt, draw=#1, line width=1pt, fill=#1!8,
            text=#1!50!black, font=\scriptsize\bfseries, align=center,
            minimum width=1.75cm, minimum height=0.72cm},
  d/.style={->, primary!45, line width=0.8pt},
  alt/.style={->, accent!70, line width=1.1pt},
  dur/.style={font=\tiny, text=neutral, inner sep=1pt}]

\node[t=primary]   (spec)    at (0,0)      {spec};
\node[t=primary]   (schema)  at (2.1,0)    {schema};
\node[t=primary]   (api)     at (4.2,0.75) {api};
\node[t=primary]   (auth)    at (4.2,-0.5) {auth};
\node[t=primary]   (ui)      at (2.1,2.0)  {ui};
\node[t=goldhl]    (harness) at (4.2,-1.85){harness};
\node[t=greenhl]   (auto)    at (6.5,-1.15){auto\_test};
\node[t=accent]    (manual)  at (6.5,0.3)  {manual\_test};
\node[t=primary]   (docs)    at (6.5,1.7)  {docs};
\node[t=primary]   (uat)     at (8.8,0.75) {uat};
\node[t=tealhl]    (rel)     at (11.0,0.75){release};

\node[dur, below=1pt of spec]    {3 d};
\node[dur, below=1pt of schema]  {2 d};
\node[dur, above=1pt of api]     {5 d};
\node[dur, below=1pt of auth]    {3 d};
\node[dur, above=1pt of ui]      {4 d};
\node[dur, below=1pt of harness] {3 d};
\node[dur, below=1pt of auto]    {2 d};
\node[dur, above=1pt of manual]  {6 d};
\node[dur, above=1pt of docs]    {2 d};
\node[dur, below=1pt of uat]     {3 d};
\node[dur, below=1pt of rel]     {1 d};

\draw[d] (spec) -- (schema);
\draw[d] (spec) -- (ui);
\draw[d] (schema) -- (api);
\draw[d] (schema) -- (auth);
\draw[d] (schema) -- (harness);
\draw[d] (api) -- (docs);
\draw[d] (api) -- (manual);
\draw[d] (auth) -- (manual);
\draw[d] (api) to[bend right=12] (auto);
\draw[d] (auth) -- (auto);
\draw[d] (harness) -- (auto);
\draw[alt] (auto) -- (uat);
\draw[alt] (manual) -- (uat);
\draw[d] (ui) to[bend left=18] (uat);
\draw[d] (docs) -- (rel);
\draw[d] (uat) -- (rel);

\node[rounded corners=3pt, fill=accent!7, draw=accent!60, line width=0.9pt,
      text=accent!55!black, font=\scriptsize, text width=5.1cm,
      align=flush left, inner sep=5pt] at (2.4,-2.55)
  {\textbf{The one choice.} \texttt{uat} needs \emph{either}
   \texttt{auto\_test} \emph{or} \texttt{manual\_test}. Automated: $3+2=5$ days.
   Manual: $6$ days. Everything else is mandatory.};

\node[font=\tiny\itshape, text=neutral, align=flush left, text width=4.0cm,
      anchor=north west] at (7.4,-2.0)
  {Mandatory work: $23$ days.\\
   Cheapest plan: $\mathbf{28}$ days.\\
   Manual route: $29$ days.\\
   $55$ reachable states.};
\end{tikzpicture}}%
{The sprint network. Ten of the eleven tasks are unavoidable; the testing step can
be reached two ways, and the cheaper route is not the shorter one in number of
tasks.}{sprint-net}

\section{Heuristics: Telling the Search Where the Goal Is}

Uniform-cost search spreads outwards in every direction equally, because it has no
idea which direction the goal is in. A \term{heuristic} gives it one.

\begin{definitionbox}[Heuristic function]
A \term{heuristic} $h(n)$ estimates the cheapest cost from the state at node $n$ to
a goal state. By convention $h(n) = 0$ for every goal state.

$h$ is \term{admissible} if it never overestimates:
\[
  0 \;\le\; h(n) \;\le\; h^{*}(n) \quad\text{for every } n,
\]
where $h^{*}(n)$ is the true cheapest completion cost from $n$.

$h$ is \term{consistent} (or \term{monotonic}) if for every node $n$ and every
successor $n'$ reached by an action of cost $c(n,n')$:
\[
  h(n) \;\le\; c(n,n') + h(n').
\]
Consistency implies admissibility. The reverse does not hold.
\end{definitionbox}

Two heuristics for the sprint, both of which a delivery manager would recognise.

\textbf{Mandatory work remaining.} Eight of the eleven tasks appear in every plan
that reaches a release: \texttt{spec}, \texttt{schema}, \texttt{api},
\texttt{auth}, \texttt{ui}, \texttt{docs}, \texttt{uat} and \texttt{release}. Sum
the durations of those not yet done. At the start this is $23$ days. It is
admissible because the tasks it counts genuinely must be done and it ignores the
testing work entirely, so it can only under-count.

\textbf{Critical path remaining.} The longest chain of dependencies still to be
traversed, taking the cheaper branch at the choice. At the start this is $16$ days.
It is admissible because every task on a chain must be done in order, so the chain
length is a lower bound on the total.

The mandatory-work estimate is much the larger of the two at the start --- $23$
against $16$ --- but, as Section~\ref{sec:dominance} shows by measurement, it is
\emph{not} larger everywhere, and that turns out to matter.

\section{Greedy Best-First Search, and Why It Is Not Enough}

Order the frontier by $h$ alone: always expand whatever \emph{looks} closest to the
goal. This is \term{greedy best-first search}, and it is fast --- on the sprint it
finds a plan after expanding ten or eleven states, against uniform-cost search's
forty-six.

It is also not optimal, and the sprint shows exactly how it fails. With the
mandatory-work heuristic, greedy search returns the \emph{manual testing} plan at
$29$ days instead of the automated one at $28$. The reason is instructive: the
harness is three days of work that reduces $h$ by nothing at all, because the
harness is not mandatory and the heuristic only counts mandatory work. Greedy search
therefore sees the harness as pure delay and refuses to build it --- which is a
recognisable human failure mode as well as an algorithmic one.

\begin{alertbox}[Greedy search ignores the road already travelled]
The defect is structural, not a matter of a better heuristic. Greedy search
minimises $h$, the estimated cost \emph{remaining}. It never consults $g$, the cost
already \emph{incurred}. So it cannot trade a more expensive first step against a
cheaper remainder --- which is precisely what investing in a test harness is.

Any decision of the form ``spend more now to spend less later'' is invisible to
greedy search. That covers tooling, automation, refactoring and nearly every
technical-debt decision in software delivery.
\end{alertbox}

\section{A* Search}

Add the two terms:
\[
  f(n) \;=\; g(n) + h(n)
\]
where $g(n)$ is the cost of the path from the start to $n$ and $h(n)$ the estimate
from $n$ to a goal. Expand the node of lowest $f$. That is \term{$A^{*}$}, and it is
the single most useful algorithm in this book.

Read $f(n)$ as \emph{the estimated cost of the cheapest complete solution that goes
through $n$}. Uniform-cost search is $A^{*}$ with $h = 0$, and greedy search is the
same frontier with $g = 0$. \dsref{fg-split} shows the two routes competing on $f$.

\dsfig{%
\begin{tikzpicture}[font=\scriptsize,
  bar/.style={draw=#1!70, fill=#1!22, line width=0.7pt},
  lbl/.style={font=\tiny, text=neutral},
  hd/.style={font=\scriptsize\bfseries, text=#1}]

% scale: 1 day = 0.34 cm
\def\s{0.34}

% --- route A: automated ------------------------------------------------
\node[hd=greenhl, anchor=west] at (0,1.55) {AUTOMATED ROUTE};
\draw[bar=primary] (0,0.85) rectangle (23*\s,1.35);
\node[font=\tiny, text=primary!60!black] at (23*\s/2,1.1)
  {mandatory work \;23 d};
\draw[bar=goldhl] (23*\s,0.85) rectangle (26*\s,1.35);
\node[lbl, anchor=south] at (24.5*\s,1.38) {harness 3};
\draw[bar=greenhl] (26*\s,0.85) rectangle (28*\s,1.35);
\node[lbl, anchor=south] at (27*\s,1.38) {auto 2};
\node[font=\scriptsize\bfseries, text=greenhl!50!black, anchor=west]
  at (28*\s+0.12,1.1) {$= 28$ d};

% --- route B: manual ---------------------------------------------------
\node[hd=accent, anchor=west] at (0,-0.35) {MANUAL ROUTE};
\draw[bar=primary] (0,-1.05) rectangle (23*\s,-0.55);
\node[font=\tiny, text=primary!60!black] at (23*\s/2,-0.8)
  {mandatory work \;23 d};
\draw[bar=accent] (23*\s,-1.05) rectangle (29*\s,-0.55);
\node[lbl, anchor=north] at (26*\s,-1.08) {manual\_test 6};
\node[font=\scriptsize\bfseries, text=accent!55!black, anchor=west]
  at (29*\s+0.12,-0.8) {$= 29$ d};

\draw[gray!45, dash pattern=on 2pt off 1.6pt] (23*\s,-1.35) -- (23*\s,1.55);

% --- the decision point -----------------------------------------------
\node[rounded corners=3pt, fill=lightbg, draw=primary!35, line width=0.9pt,
      text=primary!70!black, font=\scriptsize, text width=6.0cm,
      align=flush left, inner sep=6pt, anchor=north west] at (0,-1.75)
  {\textbf{Greedy search ($h$ only).} The harness reduces the
   mandatory-work estimate by zero, so it looks like three wasted days. Greedy
   takes the manual route: \textbf{29 days}.};

\node[rounded corners=3pt, fill=greenhl!8, draw=greenhl!50, line width=0.9pt,
      text=greenhl!45!black, font=\scriptsize, text width=6.6cm,
      align=flush left, inner sep=6pt, anchor=north west] at (6.4,-1.75)
  {\textbf{$A^{*}$ ($f = g + h$).} Both routes are carried in the frontier
   until their totals can be compared. The automated route's higher $g$ is paid
   back, and $A^{*}$ returns \textbf{28 days} --- provably the cheapest.};
\end{tikzpicture}}%
{The two routes as cost. Only an algorithm that adds the cost already spent to the
cost still expected can see that three days of harness buys four days of
testing.}{fg-split}

\subsection{Why A* Is Optimal}

\begin{worked}[The optimality argument, and where it breaks]
\textbf{Claim.} With an admissible $h$, the first goal node $A^{*}$ pops is on an
optimal path.

\smallskip
\textbf{Proof.} Let $G$ be the goal node popped, and suppose for contradiction that
its cost $g(G)$ exceeds the optimal cost $C^{*}$. Since $h(G) = 0$ for a goal,
\[
  f(G) \;=\; g(G) \;>\; C^{*}.
\]
Now consider any optimal solution path. At the moment $G$ was popped, some node $n$
on that optimal path was still in the frontier (the start is on it, and the path has
not been fully expanded, so an unexpanded node on it exists). For that $n$,
admissibility gives
\[
  f(n) \;=\; g(n) + h(n) \;\le\; g(n) + h^{*}(n) \;=\; C^{*}.
\]
So $f(n) \le C^{*} < f(G)$, and $n$ had a strictly smaller $f$ than $G$. But
$A^{*}$ pops the smallest $f$ in the frontier, so it would have popped $n$, not
$G$. Contradiction, and therefore $g(G) = C^{*}$. \hfill$\blacksquare$

\smallskip
\textbf{The step that fails when $h$ overestimates.} It is the inequality
$h(n) \le h^{*}(n)$. If $h$ overestimates at $n$, then $f(n)$ can exceed $C^{*}$;
the node on the optimal path is then buried behind $G$ in the frontier, and $A^{*}$
returns the worse answer without noticing. Concretely: give the harness a
pessimistic estimate of $10$ days remaining and $A^{*}$ reproduces greedy search's
$29$-day plan.

\smallskip
\textbf{Why graph search needs consistency as well.} The argument above assumes a
node on the optimal path is available in the frontier. With an explored set, a state
can be reached first by a dearer path, be marked explored, and then be reached again
more cheaply --- and refused. If $h$ is merely admissible, that discards the optimal
path. Consistency, $h(n) \le c(n,n') + h(n')$, guarantees $f$ never decreases along
a path, so the first time a state is popped it is already by its cheapest path. Both
heuristics in this chapter are consistent; the lab checks the weaker property,
admissibility, over all $55$ reachable states.
\end{worked}

\subsection{Dominance}
\label{sec:dominance}

If $h_{2}(n) \ge h_{1}(n)$ for every $n$, and both are admissible, then $h_{2}$
\term{dominates} $h_{1}$, and $A^{*}$ with $h_{2}$ expands no more nodes than with
$h_{1}$. The reason is direct: $A^{*}$ must expand every node with
$f(n) < C^{*}$, and a larger $h$ pushes more nodes over that threshold.

So among admissible heuristics, bigger is better --- as long as it stays admissible
and stays cheap to compute. Those are competing pressures, and the balance is a
judgement call: a heuristic that is perfect but takes as long to evaluate as solving
the subproblem has bought nothing.

Now the measurement, which is more interesting than the tidy version of this story.
At the root, mandatory work is worth $23$ and the critical path $16$, so it is
tempting to conclude that mandatory work dominates. It does not. The lab checks all
$55$ reachable states and finds $6$ on which the critical path is the \emph{larger}
of the two --- states late in the sprint where almost all the mandatory work is
finished but the testing chain still has to be walked. For instance with
\texttt{spec}, \texttt{schema}, \texttt{api}, \texttt{auth}, \texttt{docs} and
\texttt{ui} done, mandatory work is $4$ while the critical path is $9$.

So \emph{neither heuristic dominates the other}, which is the ordinary situation and
not a defect. It is also exactly the situation in which the maximum earns its keep:
$\max(h_{1},h_{2})$ is admissible, dominates both by construction, and costs one
extra comparison.

\begin{conceptbox}
Dominance is a claim about \emph{every} state, and the root is a sample of size one.
Two heuristics almost always cross somewhere, because they are usually relaxations
of different constraints and each is strongest where the constraint it kept is the
binding one.

The practical consequence: do not choose between admissible heuristics, take their
maximum. The only reason not to is evaluation cost, and a comparison is free.
\end{conceptbox}

\section{Inventing a Heuristic}

The standard technique is \term{relaxation}: delete a constraint from the problem,
solve the easier problem exactly, and use that as the estimate. It is admissible
automatically, because any solution to the real problem is also a solution to the
relaxed one, so the relaxed optimum cannot be larger.

Both heuristics here are relaxations:

\begin{itemize}[leftmargin=2.2em, itemsep=2pt, topsep=4pt]
  \item \emph{Mandatory work} relaxes away the testing requirement entirely. The
        relaxed problem --- do the eight unavoidable tasks --- is solved by adding
        up durations.
  \item \emph{Critical path} relaxes away the fact that only one task happens at a
        time, i.e.\ it assumes unlimited parallelism. The relaxed problem is the
        longest path in a directed acyclic graph.
\end{itemize}

And as Section~\ref{sec:dominance} argued, their maximum is admissible too and
dominates both. Taking the maximum of several admissible heuristics is always safe
and is the cheapest route to a stronger one.

\smallskip
\noindent\begin{longtable}{@{}L{3.9cm}L{1.5cm}L{2.1cm}L{1.9cm}L{3.9cm}@{}}
\toprule
\textbf{Algorithm} & \textbf{Cost} & \textbf{Optimal?} &
\textbf{Expanded} & \textbf{Route chosen} \\
\midrule
\endfirsthead
\toprule
\textbf{Algorithm} & \textbf{Cost} & \textbf{Optimal?} &
\textbf{Expanded} & \textbf{Route chosen} \\
\midrule
\endhead
Uniform-cost ($h=0$) & 28 & Yes & 46 & automated \\
$A^{*}$, critical path & 28 & Yes & 35 & automated \\
$A^{*}$, mandatory work & 28 & Yes & \textbf{33} & automated \\
$A^{*}$, max of the two & 28 & Yes & \textbf{33} & automated \\
Greedy, critical path & 28 & Not guaranteed & 11 & automated \\
Greedy, mandatory work & \textbf{29} & \textbf{No} & \textbf{10} & manual \\
\bottomrule
\caption{Measured on the network of \dsref{sprint-net}. $A^{*}$ keeps
uniform-cost search's guarantee while doing 28\% less work; greedy search does a
quarter of the work and, with one of the two heuristics, returns the wrong plan.}
\label{tab:informed}
\end{longtable}

The fourth row is worth a remark. Taking the maximum guarantees no \emph{more}
expansions, and here it delivers no fewer either --- $33$ both ways. The six states
where the critical path is the stronger estimate are simply not states that $A^{*}$
had to decide about. That is the normal return on a dominance improvement: a
guarantee that you will not do worse, and an empirical gain that depends on the
instance.

Read the last two rows together. Greedy search with the critical-path heuristic
happens to find the optimal plan --- and that is luck, not a property. Nothing in
the algorithm tells you which of the two rows you are in on a problem whose answer
you do not already know. That asymmetry is the whole case for $A^{*}$.

\begin{pitfall}
\begin{tight}
  \item \textbf{Forgetting the $g$ term.} Ordering the frontier by $h$ alone is
        greedy search, and it cannot recognise an investment. If your search refuses
        to build tooling, this is why.
  \item \textbf{An inadmissible heuristic.} Anything ``tuned'' upward to speed the
        search up --- multiplying a lower bound by $1.5$, say --- forfeits
        optimality. That may be a reasonable trade, but make it deliberately and
        say so.
  \item \textbf{Assuming admissibility is enough for graph search.} With an explored
        set you want consistency. Otherwise a state reached cheaply on the second
        visit is silently discarded.
  \item \textbf{A heuristic that costs more than it saves.} Expanding $30\%$ fewer
        nodes is a loss if evaluating $h$ takes twice as long. Measure time as well
        as expansions.
  \item \textbf{$h > 0$ at a goal state.} Breaks the optimality proof at its first
        line. Always check $h(\text{goal}) = 0$.
  \item \textbf{Averaging two admissible heuristics.} The average is admissible but
        is dominated by the maximum. Take $\max$, never the mean.
\end{tight}
\end{pitfall}

%---------------------------------------------------------------------
\begin{lab}[A* Over Sprint States, Two Heuristics]
Uniform-cost, $A^{*}$ and greedy search on the network of \dsref{sprint-net}, with
both heuristics, plus an exhaustive admissibility check over every reachable state.
\begin{lstlisting}[language=Python]
"""Chapter 5 lab -- informed search on the sprint network.

Reproduces Table 5.1. The headline result is the last row: greedy
search with the mandatory-work heuristic returns a 29-day plan when a
28-day plan exists, because a test harness reduces its estimate by
nothing at all.
"""

import heapq
import itertools
from collections import deque

# duration, then a conjunction of disjunctions: a task is ready when
# every group has at least one member already finished.
TASKS = {
    "spec":        (3, []),
    "schema":      (2, [{"spec"}]),
    "api":         (5, [{"schema"}]),
    "ui":          (4, [{"spec"}]),
    "auth":        (3, [{"schema"}]),
    "harness":     (3, [{"schema"}]),
    "auto_test":   (2, [{"api"}, {"auth"}, {"harness"}]),
    "manual_test": (6, [{"api"}, {"auth"}]),
    "docs":        (2, [{"api"}]),
    "uat":         (3, [{"ui"}, {"auto_test", "manual_test"}]),
    "release":     (1, [{"uat"}, {"docs"}]),
}
GOAL = "release"

# The eight tasks that appear in every plan. Everything except the two
# testing routes and the harness that one of them needs.
MANDATORY = ("spec", "schema", "api", "auth", "ui", "docs", "uat",
             "release")


def dur(t):
    return TASKS[t][0]


def ready(state):
    """Legal actions: the successor function, as a generator."""
    for t, (_, groups) in TASKS.items():
        if t not in state and all(g & state for g in groups):
            yield t


# --- the two heuristics, both relaxations --------------------------
def h_zero(state):
    return 0


def h_mandatory(state):
    """Relax away the testing requirement: just add up the work that
    cannot be avoided. 23 days at the root."""
    return sum(dur(t) for t in MANDATORY if t not in state)


def h_chain(state):
    """Relax away the one-task-at-a-time constraint: the longest
    remaining dependency chain. 16 days at the root."""
    memo = {}

    def chain(t):
        if t in state:
            return 0
        if t in memo:
            return memo[t]
        need = 0
        for g in TASKS[t][1]:
            need = max(need, min(chain(x) for x in g))
        memo[t] = need + dur(t)
        return memo[t]

    return chain(GOAL)


def h_max(state):
    """Admissible, and dominates both by construction."""
    return max(h_mandatory(state), h_chain(state))


# --- best-first search; the priority decides the algorithm ---------
def best_first(h, use_g=True):
    """use_g=True  -> A*      (f = g + h)
       use_g=False -> greedy  (f = h)

    The heap entry is (priority, counter, g, state, plan). The counter
    sits immediately after the priority so that ties are broken by
    insertion order alone. Putting g there instead would quietly let
    the cost leak into greedy search's ordering -- and greedy search
    would then find the optimal plan and the point of the lab would
    be lost.
    """
    start = frozenset()
    tie = itertools.count()
    frontier = [(h(start), next(tie), 0, start, ())]
    best = {start: 0}
    expanded = 0
    while frontier:
        _, _, g, state, plan = heapq.heappop(frontier)
        if GOAL in state:
            return g, list(plan), expanded
        if use_g and g > best.get(state, 1 << 30):
            continue                    # a cheaper path already went
        expanded += 1
        for t in ready(state):
            nxt, ng = state | {t}, g + dur(t)
            if use_g:
                if ng >= best.get(nxt, 1 << 30):
                    continue
                best[nxt] = ng
                pri = ng + h(nxt)
            else:
                if nxt in best:
                    continue
                best[nxt] = ng
                pri = h(nxt)
            heapq.heappush(frontier,
                           (pri, next(tie), ng, nxt, plan + (t,)))
    return None, None, expanded


def reachable():
    seen, q = {frozenset()}, deque([frozenset()])
    while q:
        s = q.popleft()
        for t in ready(s):
            n = s | {t}
            if n not in seen:
                seen.add(n)
                q.append(n)
    return seen


def cheapest_from(state):
    """True remaining cost h*, by uniform-cost search from state."""
    tie = itertools.count()
    fr = [(0, next(tie), state)]
    best = {state: 0}
    while fr:
        g, _, st = heapq.heappop(fr)
        if GOAL in st:
            return g
        if g > best.get(st, 1 << 30):
            continue
        for t in ready(st):
            n, ng = st | {t}, g + dur(t)
            if ng < best.get(n, 1 << 30):
                best[n] = ng
                heapq.heappush(fr, (ng, next(tie), n))
    return None


def route(plan):
    return "automated" if "auto_test" in plan else "manual"


if __name__ == "__main__":
    print("h at the root:  zero=%d  chain=%d  mandatory=%d"
          % (h_zero(frozenset()), h_chain(frozenset()),
             h_mandatory(frozenset())))
    print()
    print("algorithm                cost  expanded  route")
    print("-" * 52)
    runs = [
        ("uniform-cost (h=0)   ", h_zero, True),
        ("A* critical path     ", h_chain, True),
        ("A* mandatory work    ", h_mandatory, True),
        ("A* max of the two    ", h_max, True),
        ("greedy critical path ", h_chain, False),
        ("greedy mandatory work", h_mandatory, False),
    ]
    results = {}
    for name, h, use_g in runs:
        c, plan, exp = best_first(h, use_g)
        results[name.strip()] = (c, plan, exp)
        print("%s %5d %9d  %s" % (name, c, exp, route(plan)))

    opt = results["uniform-cost (h=0)"][0]
    print()
    print("optimal cost:", opt)

    # 1. A* is optimal with every heuristic here
    astars = ("A* critical path", "A* mandatory work",
              "A* max of the two")
    for n in astars:
        assert results[n][0] == opt, n
    # 2. and expands strictly fewer nodes than uniform-cost search
    ucs_exp = results["uniform-cost (h=0)"][2]
    for n in astars:
        assert results[n][2] < ucs_exp, n

    # 3. NEITHER heuristic dominates the other -- the interesting part
    states = reachable()
    mand_bigger = [s for s in states if h_mandatory(s) > h_chain(s)]
    chain_bigger = [s for s in states if h_chain(s) > h_mandatory(s)]
    print()
    print("reachable states: %d" % len(states))
    print("  mandatory work is the larger estimate on %d"
          % len(mand_bigger))
    print("  critical path  is the larger estimate on %d"
          % len(chain_bigger))
    assert mand_bigger and chain_bigger, "expected them to cross"
    # so max dominates both, by construction
    assert all(h_max(s) >= h_mandatory(s) for s in states)
    assert all(h_max(s) >= h_chain(s) for s in states)
    assert results["A* max of the two"][2] <= min(
        results["A* critical path"][2],
        results["A* mandatory work"][2])
    ex = min(chain_bigger, key=len)
    print("  e.g. %s: mandatory=%d, critical path=%d"
          % (" ".join(sorted(ex)), h_mandatory(ex), h_chain(ex)))

    # 4. greedy with the mandatory heuristic is NOT optimal
    assert results["greedy mandatory work"][0] > opt
    print()
    print("greedy with the mandatory heuristic returns %d, not %d"
          % (results["greedy mandatory work"][0], opt))

    # 5. all three heuristics are admissible on every reachable state
    bad = []
    for s in states:
        if GOAL in s:
            continue
        true = cheapest_from(s)
        if true is None:
            continue
        for nm, h in (("mandatory", h_mandatory), ("chain", h_chain),
                      ("max", h_max)):
            if h(s) > true:
                bad.append((nm, sorted(s), h(s), true))
    print("admissibility violations:", len(bad))
    assert not bad
    print()
    print("A* kept the guarantee and did %d%% less work than"
          % round(100 * (ucs_exp - results["A* mandatory work"][2])
                  / ucs_exp))
    print("uniform-cost search. Greedy did a quarter of the work and")
    print("got the answer wrong.")
\end{lstlisting}
The admissibility check is the part worth keeping. It is easy to invent a heuristic
that is plausible, fast and wrong, and the symptom --- a sub-optimal answer on
problems whose optimum you do not know --- is invisible in production. Checking
$h(s) \le h^{*}(s)$ exhaustively is only possible on a small instance, which is
exactly why you should do it while the instance is still small.
\end{lab}

%---------------------------------------------------------------------
\begin{checkpoint}
\begin{tightnum}
  \item The mandatory-work heuristic values the root at $23$ and the true optimum is
        $28$. Is it admissible? Would a heuristic returning $30$ at the root be
        admissible, and what would $A^{*}$ do with it?
  \item Explain in one sentence why greedy best-first search will never choose to
        build a test harness, however good the harness is.
  \item You have two admissible heuristics, $h_{1}$ and $h_{2}$, neither
        dominating the other. Give a third admissible heuristic that dominates both,
        and say why the average of the two would be a worse choice.
\end{tightnum}
\end{checkpoint}

%---------------------------------------------------------------------
\begin{chaptersummary}
\begin{tight}
  \item A \term{heuristic} $h(n)$ estimates the remaining cost. It is
        \term{admissible} if it never overestimates and \term{consistent} if
        $h(n) \le c(n,n') + h(n')$; consistency implies admissibility and is what
        graph search with an explored set needs.
  \item \term{Greedy best-first search} orders the frontier by $h$ alone. It is
        fast and not optimal, and its structural blind spot is any investment ---
        it cannot trade a dearer step now against a cheaper remainder.
  \item \term{$A^{*}$} orders by $f(n) = g(n) + h(n)$ and returns an optimal
        solution whenever $h$ is admissible. Uniform-cost search is $A^{*}$ with
        $h=0$.
  \item The optimality proof turns on $h(n) \le h^{*}(n)$: an overestimate buries
        the optimal path behind a worse goal node, and $A^{*}$ returns the worse
        answer without any sign that it has done so.
  \item Among admissible heuristics, larger is better --- $h_{2} \ge h_{1}$
        everywhere means no more expansions --- provided it stays cheap to compute.
        The maximum of several admissible heuristics is admissible and dominates
        each.
  \item \term{Relaxation} is the way to invent one: drop a constraint, solve the
        easier problem exactly. Dropping the testing requirement gives mandatory
        work; dropping the one-task-at-a-time rule gives the critical path.
  \item Measured on the sprint: $A^{*}$ matched uniform-cost search's $28$-day
        optimum while expanding $33$ states instead of $46$; greedy search expanded
        $10$ and returned a $29$-day plan.
\end{tight}
\end{chaptersummary}

\begin{reviewq}
  \item \textbf{(MCQ)} $A^{*}$ orders its frontier by: (a)~$h(n)$ (b)~$g(n)$
        (c)~$g(n)+h(n)$ (d)~depth.
  \item \textbf{(MCQ)} A heuristic is admissible if it: (a)~never underestimates
        (b)~never overestimates (c)~is always zero at the root (d)~is consistent.
  \item \textbf{(MCQ)} Uniform-cost search is $A^{*}$ with: (a)~$h=0$
        (b)~$g=0$ (c)~$h=h^{*}$ (d)~a depth limit.
  \item \textbf{(MCQ)} Given admissible $h_{1}$ and $h_{2}$, which is guaranteed
        admissible and at least as strong as both? (a)~$h_{1}+h_{2}$
        (b)~$(h_{1}+h_{2})/2$ (c)~$\max(h_{1},h_{2})$ (d)~$\min(h_{1},h_{2})$.
  \item \textbf{(Short)} State the step of the $A^{*}$ optimality proof that fails
        when $h$ overestimates, and describe the observable symptom.
  \item \textbf{(Short)} Why does graph search with an explored set want consistency
        rather than mere admissibility?
  \item \textbf{(Short)} Explain relaxation as a technique, and use it to invent an
        admissible heuristic for the problem of assigning nine developers to four
        workstreams.
  \item \textbf{(Applied)} Suppose the harness took $6$ days instead of $3$. By
        hand, compute the cost of both routes, say which is now optimal, and predict
        whether greedy search with the mandatory-work heuristic would now be
        \emph{right} by accident. Then state what that tells you about relying on
        greedy search.
\end{reviewq}
